% Part: turing-machines
% Chapter: undecidability
% Section: enumerating-tms

\documentclass[../../../include/open-logic-section]{subfiles}

\begin{document}

\olfileid[fr]{tur}{und}{enu}
\olsection{Énumération des machines de Turing}

\begin{explain}
Nous pouvons montrer que l'ensemble des machines de Turing, à renommage près, est
!!{enumerable}. Chaque machine possède en effet une description
finie. Son ensemble d'états et son alphabet de ruban sont finis.
La fonction de transition est une fonction partielle de
$Q \times \Sigma$ dans
$Q \times \Sigma \times \{\TMleft, \TMright, \TMstay\}$;
on peut donc aussi la décrire en donnant ses valeurs pour les
couples d'arguments, en nombre fini, où elle est définie.

Une subtilité demeure. Notre définition ne restreint pas les
!!{element}s qui peuvent constituer les états ou l'alphabet du
ruban : techniquement, pour chaque nombre réel, il existe par
exemple une machine qui l'emploie comme état. Pris à la lettre,
l'ensemble des machines munies de telles étiquettes arbitraires
ne serait donc pas dénombrable.\footnote{Précision éditoriale : la première affirmation source doit être comprise à renommage des états et symboles près. Les étiquettes réelles fournissent déjà une famille non dénombrable de tuples littéralement distincts.} Mais le \emph{comportement} d'une
machine ne dépend pas des objets choisis pour nommer ses états
et ses symboles. Considérons les deux machines représentées en
\olref{fig:variants}.
\begin{figure}
\begin{center}
\begin{tikzpicture}[->,>=stealth',shorten >=1pt,auto,node distance=2.8cm,
                    semithick]
  \tikzstyle{every state}=[fill=none,draw=black,text=black]

  \node[initial,state]         (A)              {$q_0$};
  \node[state]         (B) [right of=A] {$q_1$};

  \path (A) edge [bend left] node {\TMtrans{\TMstroke}{\TMstroke}{\TMright}} (B)
        (B) edge [loop above] node {\TMtrans{\TMblank}{\TMblank}{\TMright}} (B)
            edge [bend left] node {\TMtrans{\TMstroke}{\TMstroke}{\TMright}} (A);
\end{tikzpicture}\\
\begin{tikzpicture}[->,>=stealth',shorten >=1pt,auto,node distance=2.8cm,
  semithick]
\tikzstyle{every state}=[fill=none,draw=black,text=black]

\node[initial,state]         (A)              {$s$};
\node[state]         (B) [right of=A] {$h$};

\path (A) edge [bend left] node {\TMtrans{A}{A}{\TMright}} (B)
(B) edge [loop above] node {\TMtrans{\TMblank}{\TMblank}{\TMright}} (B)
edge [bend left] node {\TMtrans{A}{A}{\TMright}} (A);
\end{tikzpicture}
\end{center}
\caption{Deux variantes de la \emph{machine paire}}
\ollabel{fig:variants}
\end{figure}
Ces diagrammes définissent deux machines : $M$, dont l'alphabet
est $\Sigma = \{\TMendtape,\TMblank,\TMstroke\}$ et l'ensemble
d'états $\{q_0,q_1\}$, et $M'$, dont l'alphabet est
$\Sigma' = \{\TMendtape,\TMblank,A\}$ et les états sont
$\{s,h\}$. Leurs instructions sont par ailleurs identiques :
$M$ s'arrête sur une suite de $n$ symboles $\TMstroke$ si et
seulement si $n$ est pair, et $M'$ s'arrête sur une suite
de $n$ symboles~$A$ si et seulement si $n$ est pair.
Nous avons seulement renommé
$\TMstroke$ en~$A$, $q_0$ en~$s$ et $q_1$ en~$h$. Cet exemple
se généralise : nous pouvons identifier deux machines qui ne
diffèrent que par un tel renommage de leurs symboles et états.
Nous pouvons alors choisir des entiers positifs pour nommer
tous ces symboles et états : $1$ à la place de $\sigma_0$,
$2$ à la place de $\sigma_1$, et ainsi de suite; $\TMendtape$
devient~$1$, $\TMblank$ devient~$2$, etc. La \emph{machine paire}
prend alors la forme représentée en \olref{fig:standard-even}.
\begin{figure}
\[\begin{tikzpicture}[->,>=stealth',shorten >=1pt,auto,node distance=2.8cm,
  semithick]
\tikzstyle{every state}=[fill=none,draw=black,text=black]

\node[initial,state]         (A)              {$1$};
\node[state]         (B) [right of=A] {$2$};

\path (A) edge [bend left] node {\TMtrans{3}{3}{\TMright}} (B)
(B) edge [loop above] node {\TMtrans{2}{2}{\TMright}} (B)
edge [bend left] node {\TMtrans{3}{3}{\TMright}} (A);
\end{tikzpicture}
\]
\caption{Une \emph{machine paire} standard}
\ollabel{fig:standard-even}
\end{figure}
Nous appellerons \emph{standard} une machine dont les états et les
symboles sont des entiers positifs, et pourrons désormais nous
limiter à ces machines.\footnote{Cette appellation de « machine
standard » n'est pas consacrée par l'usage.}

Nous voulions établir que l'ensemble des machines de Turing est
!!{enumerable}; compte tenu de cette identification par renommage,
il suffit d'établir que l'ensemble des machines standard est
!!{enumerable}. Soit donc une telle
machine $M = \tuple{Q, \Sigma, q_0, \delta}$. Comment la décrire
par une suite finie d'entiers positifs ? Indiquons d'abord le
nombre d'états, puis les états eux-mêmes, le nombre de symboles,
les symboles eux-mêmes et l'état initial. Tous sont des entiers
positifs, puisque $M$ est standard. Pour décrire~$\delta$,
remarquons que son ensemble d'arguments possibles, les couples
$\tuple{q,\sigma}$, est fini parce que $Q$ et~$\Sigma$ le sont.
Les informations sur~$\delta$ se ramènent donc à la liste finie
des quintuplets
$\tuple{q, \sigma, q', \sigma', d}$ pour lesquels
$\delta(q,\sigma) = \tuple{q', \sigma', D}$, où $d$ code
la direction~$D$ : par exemple, $1$ pour~$\TMleft$,
$2$ pour~$\TMright$ et $3$ pour~$\TMstay$.

Chaque machine standard est ainsi décrite par une liste finie
d'entiers positifs, c'est-à-dire une suite $s_M \in
(\PosInt)^*$. Par exemple, la \emph{machine paire} standard est
codée par la suite
\[
2, \underbrace{1, 2}_Q, 3, \overbrace{1, 2, 3}^\Sigma, 1, \underbrace{1, 3, 2, 3, 2}_{\delta(1,3) = \tuple{2,3,R}},
\overbrace{2, 2, 2, 2, 2}^{\delta(2,2) = \tuple{2,2,R}},
\underbrace{2, 3, 1, 3, 2}_{\delta(2,3) = \tuple{1,3,R}}.
\]
\end{explain}

\begin{thm}
Il existe des fonctions de $\Nat$ vers~$\Nat$ qui ne sont
pas calculables par machine de Turing.
\end{thm}

\begin{proof}
L'ensemble $(\PosInt)^*$ des suites finies d'entiers positifs
est !!{enumerable} (\cref{sfr:siz:zigzag:prob:posint-star}).
L'ensemble des descriptions de machines standard, qui est une
partie de~$(\PosInt)^*$, est donc lui aussi dénombrable. Toute fonction
calculable par machine de Turing de $\Nat$ vers~$\Nat$ est
calculée par au moins une machine de Turing, et même par
plusieurs. En renommant ses états et ses symboles par des
entiers positifs --- notamment $\TMendtape$ par~$1$,
$\TMblank$ par~$2$ et $\TMstroke$ par~$3$ ---, nous obtenons
une machine standard qui calcule la même fonction sous ce
codage. L'ensemble des fonctions de $\Nat$ vers~$\Nat$
calculables par machine de Turing est donc dénombrable.

En revanche, l'ensemble de \emph{toutes} les fonctions de
$\Nat$ vers~$\Nat$ n'est pas !!{enumerable}
(\cref{sfr:siz:red:prob:nat-nat}). Si chacune de ces fonctions
était calculée par une machine de Turing, les descriptions
des machines qui les calculent permettraient de les énumérer
toutes. Certaines fonctions ne sont donc pas calculables par
machine de Turing.
\end{proof}

\begin{prob}
  Pouvez-vous décrire les machines de Turing sans devoir
  énumérer explicitement leurs états et les symboles de leur
  alphabet ? Vous pouvez proposer votre propre notion de machine
  « standard », mais expliquez pourquoi toute machine de Turing
  peut être simulée par une machine « standard » au sens choisi.
\end{prob}

\end{document}
